Задача сортировки "--- одна из фундаментальных задач информатики. Под сортировкой понимается
процесс упорядочивания элементов некоторого множества по заданному критерию~\cite{knuth1997art}.

Алгоритмы сортировки широко применяются в системах управления базами данных, поисковых
системах, компиляторах и других программных комплексах. Эффективность алгоритма сортировки
оказывает существенное влияние на производительность всей системы в целом.

В данной работе рассматриваются два алгоритма: \textit{быстрая сортировка} (Quicksort),
предложенная Ч.\,А.\,Р.~Хоаром в 1962 году~\cite{hoare1962quicksort}, и
\textit{сортировка слиянием} (Mergesort), основанная на принципе «разделяй и властвуй».
Оба алгоритма обладают средней временной сложностью $O(n \log n)$, однако существенно
различаются по структуре, константным коэффициентам и поведению в наихудшем случае~\cite{cormen2009algorithms}.

Целью данной работы является сравнительный анализ указанных алгоритмов: изучение их
асимптотических характеристик, реализация на языках C++ и Python, а также экспериментальное
сравнение производительности.
